
\chapter{Математическая модель квадрокоптера и разработка физически информированной нейросетевой архитектуры}
\label{ch:2}

\section{Описание платформы и ограничения навигационной системы}
\label{sec:problem}
% =========================================================

Объектом исследования является система навигации автономного квадрокоптера с максимальной взлётной
массой 1--3~кг, оснащённого тепловизионной камерой и выполняющего
термокартографирование протяжённых промышленных объектов в условиях
кратковременной потери сигнала ГНСС длительностью 3--240~с.
Предметом исследования являются методы восстановления траектории на основе
физически информированных нейронных сетей, использующих данные БИНС и
телеметрию оборотов несущих винтов, обеспечивающих адаптацию к смене
динамического режима полёта при ограниченном объёме дообучающих данных.

Квадрокоптер оснащён многоуровневой системой управления, в которой внутренние
контуры стабилизации ориентации и угловой скорости функционируют автономно от
сигнала ГНСС и способны поддерживать горизонтальный полёт при компенсации
порывов ветра.
Однако потребительские МЭМС акселерометры и гироскопы накапливают постоянно
растущую позиционную ошибку~$O(t^2)$~\cite{Chen2024,Stepanov2019}.

В данной работе рассматривается мобильный робот со следующими ограничениями
платформы. Во-первых, максимальная взлётная масса составляет 1--3~кг, что
накладывает ограничение на использование дополнительных
сенсоров~\cite{Akhmedov2022}.
Во-вторых, тепловизионная камера не может служить источником одометрии в данном
приложении, так как металлические поверхности характеризуются однородным тепловым
полем, практически лишённым угловых точек, а периодическая коррекция
неоднородности прерывает поток данных~\cite{Sokolov2023gnss}.
В-третьих, показания лидаров деградируют в запылённой и задымлённой среде
промышленных площадок~\cite{Akhmedov2022}.

Так же стоит отметить, что рассматриваемая платформа обладает нейронным
процессором~(NPU), что позволяет применять различные нейронные сети на борту
робота без повышения нагрузки на основной процессор.

Таким образом, ни один из стандартных методов навигации без ГНСС
не применим одновременно при описанных выше ограничениях.
Возникает запрос на альтернативный источник навигационной информации,
доступный на борту квадрокоптера без дополнительной нагрузки.

Альтернативным источником навигационной информации, доступным на борту без
дополнительных аппаратных средств, является телеметрия оборотов несущих
винтов; этот сигнал несёт информацию
об управляющих воздействиях, не наблюдаемую БИНС в режиме компенсации
порывов ветра~\cite{Lu2022rpm,Petrov2023nn}.

В последние годы в задачах инерциальной навигации активно развиваются гибридные архитектуры, сочетающие последовательностные модели с принципами физически информированных нейронных сетей (PINN). Такие подходы позволяют одновременно учитывать сложные временные зависимости в данных IMU и обеспечивать строгое соблюдение фундаментальных законов механики за счёт встраивания остатков дифференциальных уравнений движения в функцию потерь. Авторы отмечают, что именно это сочетание позволяет эффективно решать проблему быстрого накопления ошибок дрейфа при длительных потерях сигнала ГНСС, где традиционные методы и чисто эмпирические нейронные сети часто оказываются недостаточно эффективными.
Примером такого подхода является работа Tan и др.~\cite{Tan2023}, в которой решается задача оценки полного динамического состояния наземного транспортного средства только по данным IMU. Авторы предлагают гибридную схему, где PINN с ограничениями в виде обыкновенных дифференциальных уравнений кинематики и динамики используется для компенсации дрейфа акселерометра и гироскопа, а полученная оценка передаётся в unscented Kalman filter, работающий на многообразиях. Данный подход обеспечивает высокую точность определения трёхмерной позиции, скорости и ориентации даже при длительных сбоях GNSS. В то же время авторы отмечают в качестве недостатка необходимость использования внешнего фильтра, что усложняет полностью энд-to-энд реализацию на борту.
Аналогичный принцип гибридизации применён в работе Zhou и др.~\cite{Zhou2025}, посвящённой предсказанию раскачивания груза на мобильном кране по инерциальным измерениям. В предложенной модели LSTM-подобная последовательностная структура дополнена физическими остатками уравнений движения. Это позволило повысить физическую согласованность результатов и снизить ошибку по сравнению с обычной LSTM. Однако увеличение вычислительной сложности на длинных последовательностях ограничивает возможность применения модели в реальном времени на ресурс-ограниченных устройствах.
Для задач аэрокосмической тематики представляет интерес работа Bianchi и др.~\cite{Bianchi2024}, в которой решается задача идентификации полной динамической модели квадрокоптера, включая трансляционные и вращательные уравнения движения. Авторы используют PINN на основе многослойного персептрона с элементами последовательностной обработки, встраивая уравнения Ньютона–Эйлера непосредственно в функцию потерь. Предложенный подход продемонстрировал преимущество перед расширенным фильтром Калмана по точности при наличии неопределённостей параметров модели. Основным ограничением остаётся квадратичная вычислительная сложность относительно длины обрабатываемой последовательности, что затрудняет применение на бортовых вычислительных системах.
Задача онлайн-оценки состояния квадрокоптера в условиях ветровых возмущений рассмотрена в работе Liu и др.~\cite{Liu2025}. Авторы предложили архитектуру continual PINN с блоками, близкими по структуре к Transformer, и возможностью непрерывного дообучения. Преимуществом подхода является способность модели адаптироваться к изменяющимся внешним условиям без полной переподготовки. В то же время высокие требования к объёму памяти при обработке длинных IMU-последовательностей остаются существенным ограничением.
Лёгкие решения, ориентированные на работу в реальном времени, представлены в работе Bian и др.~\cite{Bian2025}. Авторы решают задачу оценки кинематики суставов человека по данным IMU с помощью компактной физически информированной нейронной сети. Модель сочетает последовательностную обработку с остатками уравнений движения и обеспечивает ошибку менее 5° при возможности функционирования на мобильных устройствах. Ограничением подхода является относительно небольшое число используемых IMU-датчиков, что снижает его обобщаемость на сложные многозвенные механические системы.
Для задач морской и подводной навигации представляет интерес работа An и др.~\cite{An2025}, в которой решается задача идентификации манёвров судна по инерциальным данным. Предложенная гибридная модель на основе последовательностной архитектуры с физическими ограничениями позволяет оптимизировать режим возбуждения системы для повышения точности идентификации параметров. Данный подход обеспечивает существенное снижение ошибок моделирования по сравнению с традиционными методами. Вместе с тем необходимость предварительного расчёта оптимального возбуждения усложняет практическое внедрение.
Дополнительные примеры подтверждают общую тенденцию. В работе Guang и др.~\cite{Guang2021} рассматривается прямое слияние данных IMU и GPS с помощью LSTM; авторы отмечают простоту реализации, но указывают на быстрый рост дрейфа при отсутствии физических ограничений. В исследовании Shurin и Klein~\cite{Shurin2022} предложен гибрид нейросетевой модели и уравнений движения для решения задачи dead reckoning квадрокоптера; подход демонстрирует высокую точность, однако сопровождается ростом вычислительной сложности. Обобщающий анализ представлен в обзоре Sivtsov и др.~\cite{Sivtsov2025}, где авторы констатируют преимущество гибридных моделей по обобщающей способности, но отмечают отсутствие работ, в которых линейные по сложности архитектуры SSM (Mamba) интегрированы с полноценным механизмом PINN. Наконец, в работе Fu и др.~\cite{Fu202X} показано, что гибридные PINN позволяют решать задачи предсказания динамической нагрузки с высокой интерпретируемостью, однако требуют значительного объёма обучающих данных.
Таким образом, существующие гибридные физически информированные модели на основе последовательностных архитектур демонстрируют заметный прогресс в снижении дрейфа IMU и повышении точности траекторного оценивания. В то же время в рецензируемой литературе практически отсутствуют работы, в которых последовательности с селективными пространствами состояний (архитектуры семейства Mamba/SSM) интегрированы с полноценным механизмом физически информированных нейронных сетей для задач инерциальной навигации. Ни одна из проанализированных публикаций не рассматривает комбинацию линейной по сложности Mamba-архитектуры с физическими остатками и параметрической идентификацией динамики. Указанный пробел определяет перспективное направление исследований, способное объединить высокую эффективность обработки длинных последовательностей со строгой физической согласованностью результатов.


% =========================================================
\section{Симуляционная среда с моделями датчиков}
\label{sec:simulator}
% =========================================================

Реальных открытых датасетов с синхронизированными данными МЭМС-БИНС,
ГНСС и телеметрией оборотов двигателей четырёх несущих винтов в режиме
компенсации порывов ветра в настоящее время не существует.
По этой причине эксперименты проводились в симуляторе PX4/Gazebo с плагином
RotorS~\cite{Burri2016}, являющимся стандартным бенчмарком для нейросетевой навигации БПЛА.

PX4 является открытым автопилотом с программным стеком управления, идентичным
реальным коммерческим платформам.
Gazebo обеспечивает физически корректную динамику квадрокоптера, включая аэродинамические взаимодействия несущих винтов, силу тяжести, аэродинамическое
сопротивление и модели порывов ветра.

\subsection{Симулятор Gazebo}
Симуляторы играют фундаментальную роль в современной робототехнике, позволяя генерировать большие объёмы высококачественных размеченных данных без риска повреждения реального оборудования и значительных временных затрат. Как показано в классической работе \cite{Koenig2004}, симуляция обеспечивает полную повторяемость экспериментов, точный контроль параметров окружающей среды и мгновенную доступность ground truth, что особенно важно для задач, требующих тысяч или миллионов примеров обучения. Авторы \cite{Meyer2012} подчёркивают, что такой подход существенно ускоряет процесс разработки и тестирования алгоритмов, минимизируя затраты на натурные испытания.
Gazebo представляет собой один из наиболее широко признанных открытых симуляторов, тесно интегрированный с ROS и обеспечивающий реалистичную физику, моделирование сенсоров и динамики роботов. Согласно исследованиям \cite{Borrego2018} и \cite{Chikurtev2020}, его открытый код и модульная архитектура позволяют создавать сложные сценарии с имитацией шума датчиков, ветровых возмущений и потерь сигнала ГНСС. Это делает Gazebo стандартным инструментом для подготовки данных в робототехнических исследованиях, где требуется высокая точность моделирования реальных условий.
Многочисленные публикации демонстрируют успешное использование Gazebo для сбора данных и обучения обычных нейронных сетей. Авторы \cite{Salimpour2025} показывают, как симулятор применяется для генерации датасетов в задачах навигации, управления и reinforcement learning, где обученные модели успешно переносятся на реальные платформы. Как указано в работе \cite{ZhangPINNPathPlanning}, данные из Gazebo обеспечивают необходимое разнообразие сценариев, что критично для повышения обобщающей способности нейросетей.
Для физически информированных нейронных сетей (PINN) Gazebo особенно ценен, поскольку позволяет генерировать данные с точным соблюдением физических законов движения. В публикациях \cite{Bianchi2024} и \cite{Gu2024QuadrotorPINN} подчёркивается, что симулятор обеспечивает синхронизированные данные IMU, скоростей винтов и ground truth, необходимые для построения функции потерь с остатками дифференциальных уравнений. Авторы \cite{Sahoo2025} отмечают, что такой подход минимизирует потребность в реальных измерениях и повышает физическую согласованность модели.
В настоящей работе симулятор Gazebo использован для создания синтетического датасета, включающего данные IMU, ГНСС и скоростей вращения винтов квадрокоптера Holybro X500. Это позволило эффективно обучить предложенные модели НС ПСПС и ФИ-НС ПСПС, обеспечив высокую точность восстановления траектории и последующий перенос результатов на реальное оборудование. Таким образом, применение Gazebo полностью соответствует современным практикам робототехнических исследований.

Симуляция квадрокоптера проходила в симуляторе Gazebo. Миссия полета квадрокоптера включала 500 точек. На рисунке 1 приведен скриншот программы симулятора роботов Gazebo. В качестве модели квадрокоптера был выбран стандартная для Gazebo модель Holybro X500 V2 . Характеристики квадрокоптера приведены в Приложении А. 
Подобные квадрокоптеры, такие как модель X500, широко применяются в задачах сканирования территории благодаря своей компактности, маневренности и способности обеспечивать высокоточное аэрофотосъемочное покрытие в сложных условиях рельефа \cite{Nagasawa2021, Cabreira2019}. Эти платформы позволяют проводить детальную картографию ландшафтов, мониторинг изменений поверхности и сбор данных для геоинформационных систем, минимизируя риски для операторов и повышая эффективность полевых исследований. В частности, они интегрируются с технологиями фотограмметрии и лидара для создания трехмерных моделей местности, что особенно актуально в геодезии и экологии. Такие применения подтверждаются рядом исследований, демонстрирующих практическую ценность квадрокоптеров в реальных сценариях.

На рисунке \ref{fig:x500_gazebo} представлен квадрокоптер x500 в симуляторе Gazebo

\begin{figure}[htbp]
	\centering
	\includegraphics[width=0.65\textwidth]{../src/x500_gazebo.jpg}
	\caption{Квадрокоптер Holybro X500~V2 в симуляторе Gazebo/PX4}
	\label{fig:x500_gazebo}
\end{figure}

На рисунке \ref{fig:qgc} представлен скриншот программы Qgroundcontrol с миссией для сбора датасета

\begin{figure}[htbp]
	\centering
	\includegraphics[width=0.65\textwidth]{../src/qgc.jpg}
	\caption{Миссия для обучения модели в Qgroundcontrol}
	\label{fig:qgc}
\end{figure}


Точки миссии для сбора датасета генерировались случайным образом со следующими ограничениями:
\begin{itemize}
	\item Высота над землей находится в диапазоне 5-50 метров
	\item Скорость между точками в диапазоне 1-15 м/с
	\item Расстояние между последовательными точками в диапазоне 1-10 метров.
	\item Курс квадрокоптера всегда направлен на целевую точку. 
	\item Все точки находятся в кубе с гранью длиной 150м. 
\end{itemize}
Миссия состоит из следующих этапов: взлет на высоту 50м, следование по маршрутному заданию 1000 случайно генерированных точек, остановка в последней точке, возврат к месту старта, посадка. Такая конфигурация точек позволила получить траектории с различными комбинациями значений ГНСС и БИНС. 
На рисунке 2 приведен скриншот программы QGroundControl со сгенерированной миссией 

\subsection{Модели датчиков}

В симуляторе реализованы следующие модели датчиков.
Бортовая инерциальная навигационная система включающая акселерометр моделируется с
шумом белого гауссовского шума $\sigma_a = 0{,}005$~м/с$^2/\sqrt{\text{Гц}}$,
ненулевым смещением и вибрациями от вращения роторов (частота дискретизации
200~Гц) и гироскоп $\sigma_g = 0{,}005$~рад/с$\!/\!\sqrt{\text{Гц}}$.
Эти параметры соответствуют типичным микроэлектромеханическим датчикам класса MPU-6000/ICM-42688~\cite{Matveev2020}.
Приёмник ГНСС имитируется с частотой обновления 5~Гц и позиционным шумом
$\sigma_{\text{ГНСС}} = 0{,}5$~м. Потеря сигнала моделируется обнулением выходных данных при маскировании. В задачах обучения и дообучения сигнал ГНСС является опорной траекторией
для всех рассматриваемых моделей.
Телеметрия оборотов несущих винтов представляет собой обороты всех четырёх роторов,
телеметрируемые с частотой 200~Гц непосредственно из контроллера двигателей с квантованием 10~об/мин.

Следует отметить, что перенос обученной модели на реальную аппаратную
платформу потребует верификации, так как реальные МЭМС-датчики имеют нестационарное
смещение и вибрационные помехи от моторной установки, не полностью
воспроизводимые в симуляторе~\cite{Brovko2023}.
Тем не менее использование симулятора оправдано для сравнительного анализа
архитектур, поскольку все четыре модели обучаются и тестируются в идентичных
условиях.
Эта стратегия соответствует общепринятой практике в нейросетевой
навигации~\cite{Burri2016,Kuznetsov2022}.

% =========================================================
\section{Динамическая модель квадрокоптера и информативность сигнала оборотов несущих винтов}
\label{sec:quad_model}
% =========================================================

Задача восстановления траектории при потере сигнала ГНСС сводится к оценке
вектора состояния квадрокоптера исключительно по бортовым измерениям.
В стандартном подходе таким источником информации служит БИНС, однако в режиме
компенсации порывов ветра её информативность существенно снижается.
Чтобы обосновать целесообразность включения телеметрии оборотов несущих винтов
в качестве дополнительного входного сигнала модели, необходимо проследить
причинно-следственную связь от уравнений динамики до оценки статистической
независимости сигналов.

Физическая невязка по тяге представляет собой трёхмерный вектор невязки
второго закона Ньютона~\cite{Mahony2012}:
\begin{equation}
	\boldsymbol{\varepsilon}_F = m\,\hat{\mathbf{a}} - \mathbf{F}_{\mathrm{net}}
	\in \mathbb{R}^3,
	\label{eq:resid_F}
\end{equation}
где $\hat{\mathbf{a}}$ --- предсказанное ускорение платформы,
$\mathbf{F}_{\mathrm{net}} = R_{IB}F_B + \mathbf{F}_{drag} + \mathbf{F}_{grav}$ ---
суммарная сила по трём осям инерциальной системы координат~[Н].

Суммарная сила $\mathbf{F}_{\mathrm{net}} \in \mathbb{R}^3$ в инерциальной
системе координат складывается из трёх слагаемых:
\begin{equation}
	\mathbf{F}_{\mathrm{net}} =
	\underbrace{R_{IB}\begin{bmatrix}0\\0\\F_z\end{bmatrix}}_{\text{тяга роторов}}
	+
	\underbrace{(-b_{\mathrm{drag}}\,m\,\hat{\mathbf{v}})}_{\text{аэродинамическое сопротивление}}
	+
	\underbrace{m\begin{bmatrix}0\\0\\-g_e\end{bmatrix}}_{\text{сила тяжести}},
	\label{eq:Fnet}
\end{equation}
где $F_z = k\sum_{i=1}^{4}n_i^2$ --- суммарная вертикальная тяга роторов
в связанной системе координат;
$R_{IB}$ --- матрица поворота из связанной системы координат в инерциальную,
вычисляемая из предсказанного кватерниона $\hat{q}$;
$b_{\mathrm{drag}}$ --- коэффициент аэродинамического сопротивления,
обучаемый параметр модели;
$\hat{\mathbf{v}}$ --- предсказанный вектор скорости платформы;
$m$ --- масса аппарата;
$g_e$ --- ускорение свободного падения.

Невязка по моментам:
\begin{equation}
	\boldsymbol{\varepsilon}_\tau =
	\hat{\boldsymbol{\tau}} - \mathbf{B}\,\mathbf{n}^{(\times 2)}
	\in \mathbb{R}^3.
	\label{eq:tau_df}
\end{equation}




Из уравнений~\eqref{eq:thrust}--\eqref{eq:trans_dyn} следует ключевое
наблюдение, лежащее в основе предлагаемого подхода.
При компенсации порывов ветра квадрокоптер поддерживает суммарную тягу вблизи
$F_z \approx mg$, одновременно производя быстрые асимметричные изменения
оборотов отдельных роторов для формирования управляющих моментов~\cite{Faessler2018}.
Акселерометр измеряет суммарное ускорение платформы
$\mathbf{a} = \mathbf{a}_\text{упр} + \mathbf{g}$, а гироскоп регистрирует угловую скорость корпуса $\boldsymbol{\omega}$.
Однако дифференциальные изменения тяги роторов компенсируют друг друга, а быстрые симметричные перераспределения тяги порождают угловые
возмущения, слабо различимые гироскопом на фоне шума при кратковременных порывах.
Таким образом, оба канала БИНС располагают лишь интегральной картиной воздействий
и не позволяют восстановить индивидуальный вклад каждого ротора.
Сигнал оборотов несущих винтов, напротив, непосредственно отражает
индивидуальную тягу $F_i$ по закону~\eqref{eq:thrust} и позволяет вычислить управляющие моменты по уравнению~\eqref{eq:tau_hyp}~\cite{Lu2022rpm}.
Следовательно, в режиме компенсации порывов ветра сигнал оборотов несущих
винтов обладает информацией об управляющих воздействиях, которая принципиально
недоступна ни акселерометру, ни гироскопу БИНС.
Из этого непосредственно следует, что включение телеметрии оборотов в качестве
входного сигнала нейронной сети и в качестве основы для физического
регуляризатора потерь способно повысить точность восстановления траектории
именно в режиме высокой динамики, то есть там, где задача дообучения на малых
данных наиболее сложна~\cite{Ilyasov2023pinn}.

Для количественной оценки линейной независимости сигналов БИНС и оборотов
используется коэффициент корреляции Пирсона~$r$~\cite{Kobzar2006}:
\begin{equation}
	r(X, Y) = \frac{\sum_t (X_t - \bar{X})(Y_t - \bar{Y})}
	{\sqrt{\sum_t (X_t - \bar{X})^2}\,
		\sqrt{\sum_t (Y_t - \bar{Y})^2}}.
	\label{eq:pearson}
\end{equation}
В качестве $X$ последовательно рассматриваются шесть каналов БИНС
($a_x$, $a_y$, $a_z$, $\omega_x$, $\omega_y$, $\omega_z$),
в качестве $Y$ --- дифференциальные обороты
$\Delta n_{13} = n_1 - n_3$ и $\Delta n_{24} = n_2 - n_4$,
которые непосредственно связаны с управляющими моментами по
уравнению~\eqref{eq:tau_hyp}.
Низкое значение~$|r|$ при одновременно высокой дисперсии $\Delta n_{ij}$
свидетельствует о том, что сигнал оборотов несёт информацию, линейно
независимую от БИНС, и является самостоятельным источником навигационных данных.

Для иллюстрации указанного эффекта проведён сравнительный анализ двух
характерных участков полёта продолжительностью 10~с каждый, полученных в
симуляторе PX4/Gazebo.
Выбранные участки представляют два крайних состояния динамического диапазона
платформы: спокойный режим и режим компенсации порывов ветра.
Все промежуточные режимы эксплуатации лежат в пределах диапазона, образованного
этими двумя состояниями, поэтому сопоставление именно данной пары участков
обеспечивает оценку максимально возможного контраста информативности сигналов.
Сигналы и корреляционные матрицы обоих режимов представлены на
рисунке~\ref{fig:two_regimes}, количественные характеристики сведены в
таблицу~\ref{tab:two_regimes}.

\begin{figure}[htbp]
	\centering
	\includegraphics[width=\textwidth]{../src/two_regimes_comparison.png}
	\caption{Сравнение спокойного режима и режима компенсации порывов ветра:
		9 каналов БИНС, обороты несущих винтов и корреляция между ними.\\
		\textit{Figure~1. Comparison of calm mode and wind-gust compensation mode:
			9 IMU channels, rotor RPM signals and their cross-correlation.}}
	\label{fig:two_regimes}
\end{figure}

\begin{table}[htbp]
	\centering
	\caption{Сравнение спокойного режима и режима компенсации порывов ветра\\
		\textit{Table~1. Comparison of calm flight mode and wind-gust compensation mode}}
	\label{tab:two_regimes}
	\begin{tabular}{lcc}
		\hline
		\multirow{2}{*}{Показатель} & \multicolumn{2}{c}{Режим} \\
		\cmidrule(lr){2-3}
		& Спокойный & Компенсации порывов ветра \\
		\hline
		$\sigma(\Delta n_{13})$, об/мин & 16 & 909 \\
		$\sigma(a_z)$, м/с$^2$ & 0{,}212 & 0{,}197 \\
		$\sigma(F_z)$, Н & 1{,}46 & 0{,}581 \\
		Ср. $|r|$(БИНС, $\Delta n$) & 0{,}08 & 0{,}24 \\
		\hline
	\end{tabular}
\end{table}

Анализ таблицы~\ref{tab:two_regimes} позволяет сформулировать три вывода,
непосредственно определяющих идею предлагаемого метода.


Во-первых, дисперсия вертикального ускорения БИНС~$\sigma(a_z)$ практически одинакова
в обоих режимах (0{,}212 против 0{,}197~м/с$^2$): суммарная тяга остаётся
вблизи~$mg$, и БИНС не фиксирует дифференциальных изменений тяги между
роторами, что соответствует теоретическому выводу из
уравнения~\eqref{eq:trans_dyn}.
Во-вторых, дисперсия оборотов $\sigma(\Delta n_{13})$ возрастает в режиме компенсации
порывов ветра в~57~раз.
Согласно закону~\eqref{eq:thrust}, это показывает значительную временну́ю
изменчивость управляющих сил и моментов, недоступную БИНС.
Именно эта изменчивость порождает повышенный контраст обучающего сигнала
при дообучении нейронной сети с динамическим функционалом потерь.
В третьих, средний~$|r|$ в режиме компенсации порывов ветра составляет лишь 0{,}24,
то есть остаётся низким, несмотря на 57-кратный рост дисперсии оборотов.
Это означает, что в агрессивном режиме сигнал оборотов несёт принципиально
новую, линейно независимую от БИНС информацию об управляющих воздействиях.
Следовательно, регуляризатор на основе телеметрии оборотов не дублирует
информацию, уже содержащуюся в данных БИНС, а добавляет независимый
физически обоснованный надзорный сигнал именно в тех условиях, где
переобучение на малых данных наиболее вероятно.

Таким образом, включение сигнала оборотов несущих винтов в состав входных
данных и в функционал потерь физически и информационно обосновано именно
для режима компенсации порывов ветра.
На основе изложенных выводов в следующем разделе предложена архитектура
нейронной сети с динамическим функционалом потерь, использующим
закон тяги~\eqref{eq:thrust}.




